\documentclass[12pt,a4paper]{article}

% -------------------------------------------------
% PACOTES
% -------------------------------------------------
\usepackage[utf8]{inputenc}
\usepackage[T1]{fontenc}
\usepackage[brazilian]{babel}
\usepackage{lmodern}
\usepackage{geometry}
\usepackage{setspace}
\usepackage{enumitem}
\usepackage{titlesec}
\usepackage{parskip}
\usepackage{hyperref}
\usepackage{csquotes}

% -------------------------------------------------
% REFERÊNCIAS
% -------------------------------------------------
\usepackage[
    backend=biber,
    style=abnt,
    sorting=nyt
]{biblatex}

\addbibresource{referencias.bib}

% -------------------------------------------------
% CONFIGURAÇÕES DA PÁGINA
% -------------------------------------------------
\geometry{
    top=2cm,
    bottom=2cm,
    left=2.5cm,
    right=2.5cm
}

\setstretch{1.0}

\setlength{\parindent}{0pt}
\setlength{\parskip}{4pt}

% Seções mais compactas
\titleformat{\section}
    {\large\bfseries}
    {\thesection.}
    {0.5em}
    {}

\titlespacing*{\section}
    {0pt}
    {8pt}
    {3pt}

% -------------------------------------------------
% CONFIGURAÇÃO DOS LINKS
% -------------------------------------------------
\hypersetup{
    colorlinks=true,
    urlcolor=blue,
    linkcolor=black,
    citecolor=black
}

% -------------------------------------------------
% DOCUMENTO
% -------------------------------------------------
\begin{document}

\begin{center}

{\Large \textbf{Proposta de Projeto}}\\[0.3cm]

Engenharia de Computação -- PUCPR\\[0.15cm]

\textbf{Detecção de Vulnerabilidades em Código-Fonte com Machine Learning}\\[0.15cm]

André Fabricio Wozniack -- Ricardo Hespanhol Pamplona\\[0.15cm]

Orientador: Prof. Julio Cezar Nievola

\end{center}

\vspace{0.3cm}

% -------------------------------------------------
\section{Motivação para o desenvolvimento do projeto}

Com o aumento do uso de Inteligência Artificial para gerar código,
trechos produzidos automaticamente vêm sendo incorporados a projetos
sem análise adequada de segurança, o que pode introduzir
vulnerabilidades como injeção de comandos, vazamento de dados e falhas
de gerenciamento de memória.

Estudos recentes indicam que código assistido por IA apresenta, em
média, 4,4 vezes mais vulnerabilidades que código escrito apenas por
humanos, com certos tipos de falha aparecendo com maior frequência e
templates inseguros se repetindo em diferentes repositórios, sugerindo
padrões induzidos pelos próprios modelos
\parencite{stateofcyber2026,wu2025}.

Embora ferramentas de análise estática ajudem a identificar esses
problemas, elas frequentemente geram falsos positivos e possuem
dificuldade para compreender o contexto do código, como a intenção do
desenvolvedor e chamadas entre funções
\parencite{yang2025,techscience2025}.

Por isso, há demanda por abordagens mais contextuais, baseadas em
Machine Learning e Large Language Models, que vêm sendo aplicadas à
detecção de vulnerabilidades e à priorização de alertas em código
gerado automaticamente
\parencite{wang2025,mdpi2026}.

% -------------------------------------------------
\section{Objetivo}

O objetivo do projeto é desenvolver um software capaz de analisar novos
trechos de código-fonte antes de sua integração ao sistema, identificando
padrões e erros comuns que possam resultar em vulnerabilidades de
segurança conhecidas.

A solução deverá utilizar técnicas de Machine Learning para auxiliar
na identificação e classificação dessas vulnerabilidades, buscando
complementar ferramentas tradicionais de análise estática de código.

% -------------------------------------------------
\section{Descrição sucinta das teorias e técnicas utilizadas}

O desenvolvimento do projeto envolverá conceitos de segurança de
software, análise estática de código e Machine Learning.

Serão estudadas técnicas de classificação de vulnerabilidades,
processamento e representação de código-fonte e modelos capazes de
identificar padrões associados a falhas de segurança.

Também poderão ser utilizados Large Language Models como apoio à
análise contextual do código. A avaliação da solução será realizada
por meio de métricas de classificação, como precisão, recall,
F1-score e taxa de falsos positivos, comparando os resultados obtidos
com ferramentas tradicionais de análise estática.

% -------------------------------------------------
\section{Descrição do hardware e/ou software necessários para o desenvolvimento}

\textbf{Software:}
Python, Java, Git, Docker, bancos de dados, bibliotecas de Machine
Learning e ferramentas de análise de código, como SonarQube,
linters e Dependabot.

Também serão utilizados conjuntos de dados contendo exemplos de
vulnerabilidades conhecidas para treinamento e avaliação dos modelos.

\textbf{Hardware:}
computadores pessoais dos integrantes e, caso necessário,
infraestrutura computacional local ou em nuvem para treinamento e
execução dos modelos de Machine Learning.

% -------------------------------------------------
\section{Contexto do Projeto}

O projeto será desenvolvido como Trabalho de Conclusão de Curso de
Engenharia de Computação da PUCPR.

Os integrantes serão responsáveis pelo levantamento bibliográfico,
preparação dos dados, desenvolvimento da solução, implementação dos
modelos, realização dos experimentos e avaliação dos resultados.

A solução tem como público-alvo desenvolvedores e equipes de
desenvolvimento de software que utilizam Inteligência Artificial
Generativa para auxiliar na produção de código-fonte.

% -------------------------------------------------
% REFERÊNCIAS
% -------------------------------------------------

\printbibliography[
    heading=bibintoc,
    title={Referências bibliográficas básicas}
]

% -------------------------------------------------
% ASSINATURA
% -------------------------------------------------

\vspace{1.5cm}

\begin{center}
    \rule{8cm}{0.4pt}\\
    Assinatura do orientador
\end{center}

\end{document}